% =========================================================
%  Curso de Matemática -- 9º ano do Ensino Fundamental
%  Aula 05: Introdução às Funções
% =========================================================
\documentclass[11pt,a4paper]{article}

% ---------- Idioma e codificação ----------
\usepackage[utf8]{inputenc}
\usepackage[T1]{fontenc}
\usepackage[brazil]{babel}
\usepackage{lmodern}

% ---------- Layout ----------
\usepackage[margin=2.5cm]{geometry}
\usepackage{fancyhdr}
\usepackage{enumitem}
\usepackage{multicol}
\usepackage{xcolor}
\usepackage{graphicx}

% ---------- Matemática ----------
\usepackage{amsmath, amssymb, amsthm}
\usepackage{siunitx}
\sisetup{locale = FR, group-separator = {.}}

% ---------- Figuras ----------
\usepackage{tikz}
\usetikzlibrary{calc, angles, quotes, arrows.meta}
\usepackage{pgfplots}
\pgfplotsset{compat=1.18}

% ---------- Caixas ----------
\usepackage[most]{tcolorbox}
\newtcolorbox{definicao}[1][]{colback=blue!5, colframe=blue!60!black,
  fonttitle=\bfseries, title={Definição#1}}
\newtcolorbox{exemplo}[1][]{colback=green!5, colframe=green!50!black,
  fonttitle=\bfseries, title={Exemplo#1}}
\newtcolorbox{atencao}{colback=orange!8, colframe=orange!80!black,
  fonttitle=\bfseries, title={Atenção!}}
\newtcolorbox{desafio}{colback=red!5, colframe=red!60!black,
  fonttitle=\bfseries, title={Desafio}}

% ---------- Links ----------
\usepackage{hyperref}
\hypersetup{colorlinks=true, linkcolor=blue!60!black, urlcolor=blue!60!black}

% ---------- Cabeçalho e rodapé ----------
\pagestyle{fancy}
\fancyhf{}
\lhead{Matemática -- 9º ano}
\rhead{Aula 05}
\cfoot{\thepage}

% ---------- Dados da aula ----------
\newcommand{\numaula}{05}
\newcommand{\tituloaula}{Introdução às Funções}

\begin{document}

% =========================================================
\begin{center}
  {\Large\bfseries Aula \numaula: \tituloaula}\\[4pt]
  {\large Matemática -- 9º ano do Ensino Fundamental}\\[2pt]
  Professor(a): \rule{6cm}{0.4pt} \qquad Data: \rule{3cm}{0.4pt}
\end{center}
\noindent\rule{\linewidth}{0.8pt}

\paragraph{Habilidades da BNCC:} EF09MA06.

\paragraph{Objetivos de aprendizagem:}
\begin{itemize}[nosep]
  \item Compreender função como relação de dependência entre duas variáveis.
  \item Representar funções por tabelas, leis de formação e gráficos.
  \item Identificar funções afins e quadráticas em situações do cotidiano.
\end{itemize}

\paragraph{Duração prevista:} 2 aulas de 50 minutos. % TODO: ajustar

% =========================================================
\section{Motivação}
% TODO: situação-problema ou pergunta disparadora ligada ao cotidiano.

% =========================================================
\section{A ideia de dependência entre grandezas}
% TODO: desenvolver o conteúdo.

% =========================================================
\section{Domínio, imagem e lei de formação}
% TODO: desenvolver o conteúdo.

% =========================================================
\section{O plano cartesiano e gráficos}
% TODO: desenvolver o conteúdo.

% =========================================================
\section{Função afim}
% TODO: desenvolver o conteúdo.

% =========================================================
\section{Noções de função quadrática}
% TODO: desenvolver o conteúdo.


% =========================================================
\section{Exercícios}
\begin{enumerate}[label=\textbf{\arabic*.}]
  \item % TODO: exercício de fixação
  \item % TODO: exercício de fixação
  \item % TODO: exercício contextualizado
\end{enumerate}

\begin{desafio}
  % TODO: problema desafiador (olimpíada/OBMEP).
\end{desafio}

% =========================================================
\section*{Gabarito}
% TODO: respostas comentadas.

\end{document}
